\documentclass[11pt,a4paper]{article}

% ----------------------------------------------------------------------
%  Standalone, detailed literature review for the MSc project
%  "Visual Augmentation of Audio Semantics for Accessibility"
%  Written to be read INSTEAD of the papers: each work is explained
%  (problem -> idea -> how it works -> result -> relevance to us).
% ----------------------------------------------------------------------

\usepackage[utf8]{inputenc}
\usepackage[T1]{fontenc}
\usepackage[margin=1in]{geometry}
\usepackage{booktabs}
\usepackage{longtable}
\usepackage{enumitem}
\usepackage{xcolor}
\usepackage{titlesec}
\usepackage[hidelinks]{hyperref}
\hypersetup{colorlinks=true, linkcolor=blue!50!black, urlcolor=blue!60!black, citecolor=blue!50!black}

\titleformat{\section}{\large\bfseries}{\thesection}{0.6em}{}
\titleformat{\subsection}{\normalsize\bfseries}{\thesubsection}{0.6em}{}
\setlist{nosep, leftmargin=1.4em}
\newcommand{\rel}[1]{\smallskip\noindent\textit{Relevance to this project:} #1}

\title{\textbf{A Detailed Literature Review}\\[4pt]
\large Visual Augmentation of Audio Semantics for Accessibility}
\author{Adam Gavriely \\ \small MSc Final Project --- background reading, explained}
\date{\today}

\begin{document}
\maketitle

\begin{abstract}
\noindent This document is a self-contained, explanatory literature review for the project. Its goal
is that reading it gives you the working understanding you would get from reading the primary papers:
for each area and each key work it states the \emph{problem}, the \emph{core idea}, \emph{how it
works}, the \emph{result/significance}, and \emph{why it matters for us}. The project's aim is to make
the non-speech information in a video's soundtrack --- ambient and environmental sound --- accessible
to deaf and hard-of-hearing (DHH) viewers, by detecting the sounds a hearing viewer would notice,
reasoning about which are \emph{not} already conveyed by the picture, and rendering them as
complementary visual augmentations shown alongside the video. We begin with the accessibility problem
(the ``why''), then work stage-by-stage through the technical building blocks (the ``how''), and
close by positioning the project against everything that already exists.
\end{abstract}

\tableofcontents

% ======================================================================
\section{Orientation: the problem and how the pieces fit}
% ======================================================================
The project's pipeline maps onto four mature research areas, which also organise this review:
\textbf{audio event detection} (what sounds are present, and when?), \textbf{video understanding}
(what is already visible?), \textbf{cross-modal reasoning} (which sounds are \emph{not} conveyed by
the picture, and how to depict them?), and \textbf{image/video generation} (show it). No existing
system spans all four for accessibility. The review begins with the accessibility motivation (the
``why''), then works through these building blocks (the ``how''), and closes by positioning the
project against what already exists.

% ======================================================================
\section{Accessibility: subtitles, sign language, and the case for visual augmentation}
% ======================================================================
This is the heart of the motivation, so it is treated first and in depth. The question is not merely
``can we visualise sound?'' but ``\emph{why} is a generated visual better than the accessibility tools
DHH viewers already have?'' Answering that requires understanding what subtitles and sign language do
well, where they fail for non-speech sound, and what prior work has tried.

\subsection{Subtitles and SDH: what they convey and where they break down}
Plain subtitles transcribe \emph{dialogue}. \textbf{SDH} (Subtitles for the Deaf and Hard-of-hearing)
is the richer variant that is \emph{expected} to represent non-dialogue audio too: sound effects,
music, speaker identity, manner of speaking, and audience reaction. Industry style guides such as the
\textbf{DCMP Captioning Key}~\cite{dcmp} codify the conventions: sound effects in square brackets
(\texttt{[glass shatters]}), music marked with note symbols, and --- tellingly --- the rule that the
source of a sound effect should be named \emph{``unless the source is clearly seen on screen''}, and
that ``if it can be inferred through the visuals, it is not necessary to caption it.''

That last rule is important for us: it means professional captioners already perform, by hand, exactly
the \emph{cross-modal gating} this project automates --- describe a sound only when it is not already
visible. Our Stage~5 is, in effect, an automation of an editorial principle the captioning profession
already endorses.

\medskip\noindent\textbf{Why SDH is nonetheless a poor channel for non-speech sound:}
\begin{itemize}
  \item \textbf{Low bandwidth / lost granularity.} A rich soundscape collapses to a terse label. A
    University of Sheffield study of captioned film found that in \emph{Jaws}, captions told viewers
    there was ``ominous music'' but not that the music signified the shark \emph{approaching} --- so
    DHH viewers lost the suspense the score exists to create~\cite{sheffield}. The label was present;
    the \emph{meaning} was gone.
  \item \textbf{Emotional/atmospheric loss.} An EEG study by Gerber-Mor\'on et al.~\cite{captionlimits}
    found that attending only to standard, compliant captions loses substantial emotional
    information, and that captions of non-verbal sounds tend to trigger extra \emph{attentional}
    processing rather than an emotional response --- i.e.\ the viewer reads ``[thunder]'' as a task,
    they do not \emph{feel} the storm.
  \item \textbf{Divided attention.} Captions and the picture compete for one visual channel. Eye
    tracking shows caption users fixate the text a large majority of the time and watch the scene only
    peripherally~\cite{captionvstranscript}; the effort of reading causes fatigue and missed on-screen
    content (visual gags, who is speaking). Systems like \textbf{LightSub}~\cite{lightsub} exist
    specifically to reduce this eye-movement cost. This is a warning for us: a visual augmentation is
    \emph{also} visual, so it must be designed not to worsen the very divided-attention problem it
    aims to relieve.
  \item \textbf{No spatial or temporal nuance.} Text cannot easily convey \emph{where} a sound is
    (off-screen left? approaching?) or how it evolves over time.
\end{itemize}

\subsection{Sign language: richer, but unavailable and not automatable}
Sign language can carry far more of a soundtrack than bracketed text. Skilled interpreters convey
music and sound through classifiers, mimed sound sources (an air-guitar, an instrument), sustained
holds for long notes, rhythm, and \textbf{non-manual markers} --- facial expression standing in for
pitch and intensity --- aiming to transmit the \emph{feeling} of a sound, not just its name (as seen
in theatrical and concert interpreting). So in principle sign language sets a high ceiling for
expressive, affective conveyance of non-speech sound.

Its limits for this project's setting are practical, not expressive:
\begin{itemize}
  \item \textbf{Availability.} Interpreting is live, human and costly; the overwhelming majority of
    recorded/broadcast media has no interpreter, and where one exists the focus is speech (and
    sometimes music), rarely systematic ambient/environmental sound.
  \item \textbf{Automation is unsolved.} Automatic sign-language \emph{generation} remains an open
    research problem, so it is not a deployable channel for ambient sound today.
\end{itemize}
The takeaway is that sign language shows what ``good'' looks like --- rich, atmospheric, emotional ---
but cannot be delivered automatically or at scale for the non-speech layer.

\subsection{Prior work on making sound accessible to DHH users}
Several research threads have attacked pieces of this problem; none does what we propose, but each
grounds it.
\begin{itemize}
  \item \textbf{Sound-awareness systems (real environments).} Jain et al.'s
    \textbf{SoundWatch}~\cite{soundwatch} (a smartwatch that classifies and alerts on everyday
    sounds) and \textbf{HomeSound}~\cite{homesound} (in-home displays visualising household sounds,
    deployed with DHH families) establish, from field studies, \emph{which} sounds DHH users want to
    know about --- doorbells, appliance beeps, alarms, a baby crying, a dog, running water, sirens,
    a running vehicle. \textbf{ProtoSound}~\cite{protosound} lets users personalise a sound
    recogniser with a few examples. Key lesson for us: these target \emph{alerting/functional} sounds
    in the user's own environment; \emph{media} accessibility for \emph{atmospheric/narrative} sound
    is a different, comparatively under-served need. But the sound taxonomy transfers.
  \item \textbf{Enriching captions (media).} A line of HCI work augments \emph{text} with visual cues:
    visualising speech prosody and emotion in captions~\cite{prosodycaptions}, visualising speech
    styles~\cite{speechstyles}, and generatively adapting non-speech captions
    (\textbf{CapTune}~\cite{captune}). These show the community sees the same gap we do --- text alone
    is insufficient --- but they stay within the caption/text paradigm and focus on \emph{speech}
    paralinguistics, not generated imagery for ambient sound.
\end{itemize}

\subsection{The case for visual augmentation --- a three-way comparison}
Putting it together, the three channels have complementary strengths:

\begin{center}
\begin{longtable}{@{}p{3.2cm}p{3.5cm}p{3.5cm}p{3.5cm}@{}}
\toprule
& \textbf{Subtitles / SDH} & \textbf{Sign language} & \textbf{Visual augmentation (this project)} \\
\midrule
\endhead
Speech & excellent & excellent & not the target (handled by captions) \\
Non-speech ambient sound & terse label only & rich, if present & rich, generated depiction \\
Emotion / atmosphere & largely lost & strong & potentially strong (imagery) \\
Spatial / ``where'' info & none & possible (space/classifiers) & possible (directional cue) \\
Availability / scale & ubiquitous, cheap & rare, costly & automatable \\
Automatable today & yes & no (generation unsolved) & yes (off-the-shelf models) \\
Added cognitive load & competes with scene & separate focus point & must be designed carefully \\
\bottomrule
\end{longtable}
\end{center}

\noindent Visual augmentation occupies the sweet spot: it is higher-bandwidth than an SDH label,
automatable unlike interpreting, and it targets exactly the non-speech layer both other channels
under-serve. It does not replace subtitles or sign language --- it complements them, adding the
missing ``what does it look like / where is it / what mood does it set'' that a label cannot carry.
Two design cautions follow directly from the literature above: (1) because it is visual, it must not
worsen divided attention --- argue for \emph{few, well-chosen} augmentations (aggressive gating) and a
stable, peripheral placement; (2) atmospheric/emotional fidelity, not just correct labels, is what
actually helps DHH viewers, which is why generation (imagery) beats classification (a word).

\rel{This is the project's justification. The novelty is not ``visualise sound'' (others visualise
sound) but ``\emph{decide} which sounds are unconveyed by the picture and depict \emph{those}, for
media accessibility'' --- a combination no prior system provides.}

% ======================================================================
\section{Audio Event Detection: turning a soundtrack into labelled events}
% ======================================================================
This is the engine that answers ``what sounds are present, and when?''. The field is organised around
one dataset and two model families.

\subsection{AudioSet --- the ontology everything is built on}
\textbf{AudioSet}~\cite{audioset} (Gemmeke et al., Google, 2017) is a human-labelled dataset of
$\sim$2\,million 10-second YouTube clips organised under an ontology of \textbf{527 sound classes}
(speech, animals, vehicles, water, alarms, music, and so on). It is the de-facto vocabulary and
training set for modern audio recognition. \emph{Its strength and weakness for us:} broad coverage of
everyday sounds (water, cat, dog, rain, sirens are all present), but a large fraction of the 527 are
music/instrument classes, granularity is uneven, and it is a fixed, closed set.

\subsection{PANNs --- the practical workhorse (what we use)}
\textbf{PANNs}~\cite{panns} (``Pretrained Audio Neural Networks'', Kong et al., 2020) are CNNs trained
on AudioSet, the best-known being \textbf{CNN14}. \emph{How they work:} the audio is turned into a
log-mel spectrogram (a time--frequency image), and a CNN scores each of the 527 classes. The clip-level
model reports what sounds occur; the \emph{DecisionLevelMax} variant produces \emph{frame-level}
outputs --- a probability per class at each $\sim$10\,ms frame --- giving time boundaries. PANNs reached
a then-state-of-the-art mean average precision of \textbf{0.439} on AudioSet tagging. They are light,
fast (sub-second on CPU for a short clip), and permissively usable.
\rel{PANNs CNN14 (frame-level) is our implemented Stage-4 detector: it gives (label, start, end,
confidence) events plus a timeline. It is the recommended first pass; its closed-set/granularity
limits (e.g.\ it under-detects a regional air-raid siren) motivate an open-vocabulary companion.}

\subsection{AST and BEATs --- transformer detectors (the quality path)}
The \textbf{Audio Spectrogram Transformer (AST)}~\cite{ast} (Gong et al., 2021) replaced the CNN with a
pure transformer over spectrogram patches, capturing global context better. \textbf{BEATs}~\cite{beats}
(Chen et al., 2023) is the current practical state of the art: it is pre-trained with an
\emph{acoustic tokenizer} (a self-supervised scheme that learns discrete audio tokens, analogous to
how language models are pre-trained), yielding strong general audio representations. BEATs is the
baseline encoder in the DCASE sound-event-detection challenges, and appears again as the audio encoder
inside audio captioners and even the TempoTokens video generator.
\rel{BEATs is our planned Stage-4 quality upgrade from PANNs, and appears across several other systems
we may use --- a good encoder to standardise on.}

% ======================================================================
\section{Audio Captioning: describing a soundscape in words}
% ======================================================================
Where tagging outputs labels, \emph{automated audio captioning} (AAC) outputs a sentence: ``a train
approaches as birds chirp in the background.'' This is a richer, compositional representation and a
natural bridge to an image-generation prompt.

\textbf{Datasets.} \textbf{AudioCaps}~\cite{audiocaps} ($\sim$46\,k AudioSet clips with human
captions) and \textbf{Clotho}~\cite{clotho} are the standard benchmarks.

\textbf{CLAP} (Contrastive Language--Audio Pretraining)~\cite{clap} is the workhorse underneath much of
this: like CLIP for images, it learns a shared embedding space for audio and text, so you can measure
how well a sound matches a text description --- \emph{including arbitrary descriptions it was never
explicitly trained to classify}. This makes CLAP an \textbf{open-vocabulary} tool.

\textbf{Modern captioners} pair a strong audio encoder with a language decoder:
\textbf{EnCLAP}~\cite{enclap} (neural audio codec + CLAP + a language model) and
\textbf{SLAM-AAC}~\cite{slamaac} (current state of the art on Clotho/AudioCaps, using paraphrase
augmentation and CLAP-based re-ranking of candidate captions).
\rel{Captioning gives us (a) a natural way to turn a detected sound into a generation prompt, (b) the
required ``audio captioning only'' \emph{baseline}, and (c) via CLAP, both an open-vocabulary
detection option and an audio$\leftrightarrow$image relevance metric for evaluation.}

% ======================================================================
\section{Audio Understanding: large audio-language models (LALMs)}
% ======================================================================
LALMs fold ASR, sound-event reasoning and captioning into one instructable model: an audio encoder
$+$ a connector $+$ a large language model, so you can \emph{ask} about a sound in natural language.
\begin{itemize}
  \item \textbf{Qwen-Audio / Qwen2-Audio}~\cite{qwenaudio} --- strong general audio understanding;
    can be prompted both to transcribe speech and to describe non-speech sounds; outperformed prior
    methods across many audio benchmarks.
  \item \textbf{SALMONN}~\cite{salmonn} --- combines speech, audio-event and music encoders on a
    frozen LLM, giving general ``hearing'' abilities (ASR, translation, audio captioning) in one
    model.
  \item \textbf{Audio Flamingo 3}~\cite{audioflamingo3} --- a fully open, long-context LALM supporting
    multi-turn, multi-audio dialogue.
  \item \textbf{MMAU}~\cite{mmau} --- a benchmark for how well such models actually \emph{reason}
    about audio, useful for justifying model choices.
\end{itemize}
\rel{A LALM could collapse our Stages 3--4 (ASR $+$ event detection) into one component, and offers an
\emph{open-vocabulary} route (describe any sound in text). Trade-off: heavier and less individually
inspectable than PANNs $+$ Whisper as separate, swappable modules. Plan: keep the modular path as the
default; treat a LALM (Qwen2-Audio) as an alternative backend and an ablation.}

% ======================================================================
\section{Vision-Language and Multimodal LLMs: understanding the picture}
% ======================================================================
Two of our stages need to ``see'': Stage~2 (what is already visible?) and Stage~7 (describe the
augmented scene for evaluation).
\begin{itemize}
  \item \textbf{Qwen2.5-VL}~\cite{qwenvl} --- the leading open option: native dynamic-resolution
    vision, explicit \emph{video} understanding with temporal grounding (it can point to \emph{when}
    an event happens), released at 3B/7B/72B. The 7B fits a mid-range GPU; the 3B is a CPU/low-VRAM
    fallback.
  \item \textbf{InternVL3}~\cite{internvl} and the \textbf{LLaVA}~\cite{llava} family are the main
    open alternatives; LLaVA (``visual instruction tuning'') popularised the connector-\,+\,-frozen-LLM
    recipe most VLMs now follow.
\end{itemize}
\rel{Qwen2.5-VL-7B is our candidate for Stage~2 (scene/visible-object description) and doubles as the
Stage-7 \emph{describer}. Because the same family spans understanding \emph{and} evaluation, we must
use different instances/prompts for generation vs.\ judging so a model never grades its own output.}

% ======================================================================
\section{Sound-Source Localization: is the source on-screen or off?}
% ======================================================================
This area is the technical key to our core ``gap-aware'' decision, and is easy to overlook. The
question ``is the thing making this sound actually visible in the frame?'' has its own literature.
\begin{itemize}
  \item \textbf{Sound-source localization via cross-modal alignment}~\cite{ssloc} (Senocak et al.,
    ICCV 2023) learns to point to \emph{where} in an image a sound originates by aligning audio and
    visual features --- and, by extension, whether such a location exists at all.
  \item \textbf{Object-aware localization}~\cite{objaware} grounds the sound in specific visible
    objects via audio-visual scene understanding.
  \item \textbf{Off-screen sound separation}~\cite{offscreensep} explicitly models the case where a
    sound has \emph{no} on-screen source, estimating an on-screen-vs-off-screen probability and
    separating off-screen sound (sometimes using binaural cues for direction).
\end{itemize}
\rel{These give a \emph{quantitative} ``is the source visible?'' signal --- exactly our gate ---
without relying on an LLM's judgement. Recommended design: use a localization signal to \emph{ground}
the LLM's gating decision, and offer the ablation ``LLM-only vs.\ LLM$+$localization''.}

% ======================================================================
\section{Multimodal Reasoning and its Failure Modes}
% ======================================================================
Stage~5 asks a model to fuse scene, speech and audio events and decide what to depict. The literature
warns this is where things break:
\begin{itemize}
  \item Multimodal LLMs \textbf{hallucinate} --- inventing content not in the input; ``more thinking,
    less seeing'' can amplify it. \textbf{Fork-Merge decoding}~\cite{forkmerge} mitigates this in
    audio-visual LLMs by reasoning over each modality separately before merging.
  \item Models suffer \textbf{cross-modal attention imbalance} --- over-attending to one modality;
    ``Pay More Attention To Audio''~\cite{avattn} documents and corrects this in audio-language
    models.
\end{itemize}
\rel{Direct implications for Stage~5: feed the LLM \emph{structured} evidence (the visible-object
list, the transcript, the event labels) rather than raw media; force structured JSON output with a
required ``reason'' field; and forbid it from inventing sounds/objects not in its inputs.}

% ======================================================================
\section{Audio-to-Image Generation: showing a sound}
% ======================================================================
Two philosophies, both relevant.

\textbf{(1) Direct audio-to-image.} \textbf{Sound2Scene}~\cite{sound2scene} (Sung-Bin et al., CVPR
2023) learns, from unlabelled video alone, to map audio into the latent space of a pretrained image
generator and synthesise a matching scene directly from sound.
\textbf{SonicDiffusion}~\cite{sonicdiffusion} conditions a pretrained diffusion model on audio. These
are ``sound in, picture out'' with \emph{no} language step and \emph{no} reasoning about the existing
video --- which makes them the natural \emph{direct-audio-to-image baseline} the proposal requires.

\textbf{(2) Text-conditioned generation.} \textbf{Stable Diffusion XL (SDXL)}~\cite{sdxl} and
\textbf{FLUX.1}~\cite{flux} are the open state of the art in text-to-image. Our method sits here, but
with the crucial twist that the \emph{prompt is chosen by cross-modal reasoning} (Stage~5), not
produced blindly from the audio.
\rel{Generate augmentations with SDXL/FLUX from a Stage-5 prompt; use Sound2Scene as the
``blind audio$\to$image, no visual reasoning'' baseline. The contrast --- reasoned, gap-aware prompting
vs.\ blind generation --- is exactly what the evaluation isolates.}

% ======================================================================
\section{Audio-to-Video Generation: TempoTokens (in depth)}
% ======================================================================
\textbf{TempoTokens}~\cite{tempotokens} (Yariv, Gat, Benaim, Wolf, Schwartz, Adi; AAAI 2024) is
directly relevant and worth understanding fully.

\textbf{Problem.} Generate a short video that is aligned to an input audio clip both \emph{globally}
(the whole video matches the sound's semantics) and \emph{temporally} (each segment of audio maps to
the corresponding segment of video) --- for diverse, natural sounds, not just talking heads.

\textbf{How it works.} It reuses two frozen pretrained models and trains only a small bridge between
them: (i) a frozen \textbf{BEATs} audio encoder produces embeddings; (ii) a lightweight
\textbf{adapter} (``AudioMapper'', an MLP) maps those embeddings --- using windows of varying sizes to
capture local-to-global context --- into \emph{pseudo-text-tokens}, i.e.\ vectors in the same space the
text-to-video model expects from a text prompt; (iii) a frozen \textbf{ModelScope} latent-diffusion
text-to-video model turns those tokens into video. Because the audio is converted into the model's
``text token'' language, the system can be conditioned on audio, text, or \emph{both} (a first).

\textbf{Evaluation and results.} Trained/evaluated on \textbf{Landscape}, \textbf{AudioSet-Drums} and
\textbf{VGGSound}. The authors introduce \textbf{AV-Align}, a metric that detects energy peaks in the
audio and in the generated video and measures how well they line up in time --- a way to quantify
temporal synchrony. TempoTokens reports better audio alignment, visual quality and diversity than
prior audio-to-video methods, plus a human study.

\textbf{Limitations.} It is bounded by the frozen T2V model's capability; temporal conditioning is
capped at 24 frames (short clips); and audio--visual discrepancies can arise (the sound may not fully
determine the scene). It \emph{generates/replaces} footage and never analyses an existing video ---
i.e.\ it does not do cross-modal gating.
\rel{Three concrete uses: (1) the \emph{direct audio-to-video baseline} (the video analogue of
Sound2Scene, contrasting with our gap-aware method) --- the primary, low-effort use; (2) an
off-the-shelf motion-generation backend for our stretch goal; (3) its \textbf{AV-Align} metric for
scoring temporal alignment of moving augmentations. It shares BEATs and VGGSound with us. It
complements, not replaces, our contribution. (Co-authored by I.\ Schwartz.)}

% ======================================================================
\section{Datasets}
% ======================================================================
\begin{itemize}
  \item \textbf{AudioSet}~\cite{audioset} --- 2\,M clips, 527-class ontology; the training basis for
    PANNs/BEATs and our detection vocabulary.
  \item \textbf{AudioCaps}~\cite{audiocaps} / \textbf{Clotho}~\cite{clotho} --- audio captioning
    benchmarks.
  \item \textbf{VGGSound}~\cite{vggsound} --- $\sim$200\,k 10-second YouTube clips, 309 classes, each
    with a \emph{visible} sound source. Clean single-event ambient clips; its ``source on screen''
    property can be inverted to build ``source off-screen'' test cases that stress our gate. Also used
    by TempoTokens.
  \item \textbf{UnAV-100}~\cite{unav100} --- $\sim$10.8\,k untrimmed videos, 30\,k$+$ audio-visual
    events, 100 classes, avg 2.8 (often overlapping) events per video; CC-BY 4.0. Ideal for
    \emph{mixed} scenes with concurrent sounds.
  \item \textbf{MovieNet}~\cite{movienet} --- 1{,}100 movies with rich annotations (42\,k scene
    boundaries, character/place/action/style tags); best source for cinematic \emph{acoustic-event}
    scenes with narrative context.
  \item \textbf{Landscape} / \textbf{AudioSet-Drums} --- used by TempoTokens; narrower domains.
\end{itemize}
\rel{Our $\sim$300-clip benchmark samples VGGSound (clean ambient/off-screen cases), UnAV-100
(mixed/overlapping), and MovieNet (narrative acoustic events), across the three proposal scenarios.}

% ======================================================================
\section{Evaluation Methods}
% ======================================================================
\begin{itemize}
  \item \textbf{CLIP}~\cite{clip} --- shared image--text embedding; measures how well a generated image
    matches a text description (relevance).
  \item \textbf{CLAP}~\cite{clap} --- the audio analogue; measures audio$\leftrightarrow$text (and,
    via a shared space, audio$\leftrightarrow$image) relevance.
  \item \textbf{LLM-as-judge} --- an independent LLM scores semantic consistency between the
    augmentation's described content and a reference derived from the original multimodal input; the
    proposal's automatic protocol (VLM-describe $\rightarrow$ LLM-judge). Must be validated against a
    small human-labelled subset and must not grade its own model's output.
  \item \textbf{AV-Align}~\cite{tempotokens} --- temporal alignment of audio and (generated) video via
    energy-peak matching; for evaluating \emph{moving} augmentations.
  \item \textbf{Gating accuracy (ours)} --- precision/recall of the ``augment vs.\ stay silent''
    decision on a hand-labelled subset; the metric that most directly measures the core contribution.
\end{itemize}

% ======================================================================
\section{Synthesis: where this project sits}
% ======================================================================
Reading across the areas, the landscape is:
\begin{itemize}
  \item \textbf{Accessibility} research has thoroughly documented that non-speech sound is poorly
    served by subtitles, that sign language is expressive but unavailable/unautomatable, and has built
    sound-\emph{awareness} tools for real environments and caption-\emph{enrichment} tools for speech
    --- but not generated imagery for ambient sound in media, gated by what is already visible.
  \item \textbf{Detection, captioning, LALMs, VLMs, localization} are all mature enough to assemble
    off-the-shelf; the pieces exist.
  \item \textbf{Generation} (audio-to-image, audio-to-video) exists but in a \emph{blind} form ---
    ``sound in, picture out'' --- with no reasoning about the existing video. Those systems
    (Sound2Scene, TempoTokens) are precisely our \emph{baselines}.
\end{itemize}
The gap nobody fills is the combination: \emph{detect non-speech sound, reason about the audio--visual
gap (what is heard but not seen), and generate a complementary visual for accessibility.} That
combination --- task formulation, cross-modal gating, benchmark, and an evaluation protocol for it ---
is the contribution. Everything reviewed here is either a component we glue together or a baseline we
measure against.

% ======================================================================
\begin{thebibliography}{99}
% ======================================================================
\footnotesize
\bibitem{audioset} J.\ F.\ Gemmeke et al., ``AudioSet: An ontology and human-labeled dataset for audio events,'' \emph{ICASSP}, 2017.
\bibitem{panns} Q.\ Kong et al., ``PANNs: Large-Scale Pretrained Audio Neural Networks for Audio Pattern Recognition,'' \emph{IEEE/ACM TASLP}, 2020. arXiv:1912.10211.
\bibitem{ast} Y.\ Gong, Y.-A.\ Chung, J.\ Glass, ``AST: Audio Spectrogram Transformer,'' \emph{Interspeech}, 2021. arXiv:2104.01778.
\bibitem{beats} S.\ Chen et al., ``BEATs: Audio Pre-Training with Acoustic Tokenizers,'' \emph{ICML}, 2023. arXiv:2212.09058.
\bibitem{audiocaps} C.\ D.\ Kim et al., ``AudioCaps: Generating Captions for Audios in The Wild,'' \emph{NAACL}, 2019.
\bibitem{clotho} K.\ Drossos et al., ``Clotho: An Audio Captioning Dataset,'' \emph{ICASSP}, 2020.
\bibitem{clap} B.\ Elizalde et al., ``CLAP: Learning Audio Concepts from Natural Language Supervision,'' \emph{ICASSP}, 2023.
\bibitem{enclap} J.\ Kim et al., ``EnCLAP: Combining Neural Audio Codec and Audio-Text Joint Embedding for Automated Audio Captioning,'' \emph{ICASSP}, 2024. [verify authors]
\bibitem{slamaac} W.\ Chen et al., ``SLAM-AAC: Enhancing Audio Captioning with Paraphrasing Augmentation and CLAP-Refine through LLMs,'' 2024. arXiv:2410.09503.
\bibitem{qwenaudio} Y.\ Chu et al., ``Qwen-Audio / Qwen2-Audio: Advancing Universal Audio Understanding,'' 2023--2024. arXiv:2311.07919.
\bibitem{salmonn} C.\ Tang et al., ``SALMONN: Towards Generic Hearing Abilities for Large Language Models,'' \emph{ICLR}, 2024.
\bibitem{audioflamingo3} ``Audio Flamingo 3: Advancing Audio Intelligence,'' 2025. arXiv:2507.08128. [verify authors]
\bibitem{mmau} S.\ Sakshi et al., ``MMAU: A Massive Multi-Task Audio Understanding and Reasoning Benchmark,'' 2024. arXiv:2410.19168.
\bibitem{qwenvl} Qwen Team, ``Qwen2.5-VL Technical Report,'' 2025. arXiv:2502.13923.
\bibitem{internvl} ``InternVL3: Native Multimodal Pretraining,'' 2025. [verify authors/venue]
\bibitem{llava} H.\ Liu et al., ``Visual Instruction Tuning (LLaVA),'' \emph{NeurIPS}, 2023.
\bibitem{ssloc} A.\ Senocak et al., ``Sound Source Localization is All about Cross-Modal Alignment,'' \emph{ICCV}, 2023. arXiv:2309.10724.
\bibitem{objaware} ``Object-aware Sound Source Localization via Audio-Visual Scene Understanding,'' 2025. arXiv:2506.18557. [verify authors]
\bibitem{offscreensep} ``Off-Screen Sound Separation Based on Audio-visual Pre-training Using Binaural Audio,'' 2023. [verify authors]
\bibitem{forkmerge} ``Fork-Merge Decoding: Enhancing Multimodal Understanding in Audio-Visual Large Language Models,'' 2025. arXiv:2505.20873. [verify authors]
\bibitem{avattn} ``Pay More Attention To Audio: Mitigating Imbalance of Cross-Modal Attention in Large Audio Language Models,'' 2025. arXiv:2509.18816. [verify authors]
\bibitem{sound2scene} K.\ Sung-Bin et al., ``Sound to Visual Scene Generation by Audio-to-Visual Latent Alignment (Sound2Scene),'' \emph{CVPR}, 2023. arXiv:2306.11504.
\bibitem{sonicdiffusion} ``SonicDiffusion: Audio-Driven Image Generation and Editing with Pretrained Diffusion Models,'' 2024. [verify authors]
\bibitem{sdxl} D.\ Podell et al., ``SDXL: Improving Latent Diffusion Models for High-Resolution Image Synthesis,'' 2023. arXiv:2307.01952.
\bibitem{flux} Black Forest Labs, ``FLUX.1,'' 2024. \url{https://blackforestlabs.ai/}.
\bibitem{tempotokens} G.\ Yariv, I.\ Gat, S.\ Benaim, L.\ Wolf, I.\ Schwartz, Y.\ Adi, ``Diverse and Aligned Audio-to-Video Generation via Text-to-Video Model Adaptation (TempoTokens),'' \emph{AAAI}, 2024. arXiv:2309.16429. \url{https://arxiv.org/abs/2309.16429}
\bibitem{dcmp} Described and Captioned Media Program (DCMP), ``Captioning Key --- Guidelines and Preferred Techniques'' (Sound Effects, Music, Non-Speech Information). \url{https://dcmp.org/learn/captioningkey}
\bibitem{captionlimits} V.\ Gerber-Mor\'on et al., ``Limitations of Standard Accessible Captioning of Sounds and Music for DHH People: An EEG Study,'' \emph{Frontiers in Integrative Neuroscience}, 2020. [verify authors]
\bibitem{sheffield} University of Sheffield / SubTxt Creative, study on captions and suspense in film and TV, 2024. [verify primary publication]
\bibitem{captionvstranscript} R.\ Kushalnagar et al., ``Captions versus Transcripts for Online Video Content,'' \emph{W4A}, 2013. [verify authors]
\bibitem{lightsub} ``LightSub: Unobtrusive Subtitles with Reduced Information and Decreased Eye Movement,'' \emph{MTI}, 2024. [verify authors]
\bibitem{prosodycaptions} ``Visualization of Speech Prosody and Emotion in Captions: Accessibility for Deaf and Hard-of-Hearing Users,'' \emph{CHI}, 2023.
\bibitem{speechstyles} ``Visualizing Speech Styles in Captions for Deaf and Hard-of-Hearing Viewers,'' \emph{Int.\ J.\ Human-Computer Studies}, 2024. [verify authors]
\bibitem{captune} ``CapTune: Adapting Non-Speech Captions With Anchored Generative Models,'' 2025. arXiv:2508.19971. [verify authors]
\bibitem{soundwatch} D.\ Jain et al., ``SoundWatch: Smartwatch-based Deep Learning Approaches to Support Sound Awareness for DHH Users,'' \emph{ASSETS}, 2020.
\bibitem{homesound} D.\ Jain et al., ``HomeSound: An Iterative Field Deployment of an In-Home Sound Awareness System for DHH Users,'' \emph{CHI}, 2020.
\bibitem{protosound} D.\ Jain et al., ``ProtoSound: A Personalized and Scalable Sound Recognition System for DHH Users,'' \emph{CHI}, 2022. arXiv:2202.11134.
\bibitem{vggsound} H.\ Chen et al., ``VGGSound: A Large-scale Audio-Visual Dataset,'' \emph{ICASSP}, 2020. arXiv:2004.14368.
\bibitem{unav100} T.\ Geng et al., ``Dense-Localizing Audio-Visual Events in Untrimmed Videos (UnAV-100),'' \emph{CVPR}, 2023. arXiv:2303.12930.
\bibitem{movienet} Q.\ Huang et al., ``MovieNet: A Holistic Dataset for Movie Understanding,'' \emph{ECCV}, 2020. arXiv:2007.10937.
\bibitem{clip} A.\ Radford et al., ``Learning Transferable Visual Models From Natural Language Supervision (CLIP),'' \emph{ICML}, 2021.
\end{thebibliography}

\end{document}
